\section{What We Formalized}\label{sec:corpus}

Table~\ref{tab:corpus} lists the sixteen sources. For each it gives the
object, the method of the completeness proof, the size of our
development, and the form of the theorem we proved. A
\emph{presentation} means that the relations present the whole matrix
group; a \emph{sub-presentation} means that they present the subgroup,
or submonoid, that the generators generate, which for circuits is the
group the paper is about. ``Every $n$'' means every number of wires or
every dimension. The line counts are non-blank lines of Agda. Every
presentation theorem has the form of Section~\ref{sec:theorem-form},
the qudit Clifford ones with a model in place of the matrix group.

\begin{table}
\caption{The sources and our developments. Modules and lines are those
of the development's own folder, which includes the modules of the
matrix layer that first appeared there; the library proper is not
counted. Measured on 22 October 2026. The Hadamard-$\Pi$ theorem takes
its hard case from the row below it (Section~\ref{sec:gap}).}
\label{tab:corpus}
\footnotesize
\setlength{\tabcolsep}{4pt}
\begin{tabularx}{\linewidth}{@{}l Y l r r l@{}}
\toprule
Source & Object & Method & Mod. & Lines & Result \\
\midrule
\multicolumn{6}{@{}l}{\emph{Groups of two-level matrices}}\\
\citet{greylyn2014} & $U_4(\Z[\frac1{\sqrt2},i])$ & synthesis & 31 & 8{,}490 & presentation, $n\le4$ \\
\citet{bian2021un} & $U_n(\Z[\frac12,i])$ & synthesis & 37 & 9{,}006 & presentation, every $n$ \\
\citet{li2021on} & $O_n(\Z[\frac12])$, $L_n$ & synthesis & 68 & 15{,}528 & presentations, every $n$ \\
\citet{fang2026hadamard} & $O_n(\Z[\frac1{\sqrt2}])$ & synthesis & 44 & 11{,}205 & presentation, every $n$ \\
\citet{clement2019thesis} & $O_n(\Z[\frac1{\sqrt2}])$ & synthesis & 40 & 10{,}427 & presentation, every $n$ \\
\midrule
\multicolumn{6}{@{}l}{\emph{Circuits}}\\
\citet{selinger2015clifford} & Clifford & normal forms & 31 & 11{,}555 & sub-presentation, every $n$ \\
\citet{bian2026qupit} & Clifford, odd prime $p$ & coset tower & 215 & 70{,}088 & presentation of a model, every $n$ \\
\citet{makary2021real} & real stabilizer & normal forms & 30 & 12{,}010 & sub-presentation, every $n$ \\
\citet{amy2018cnotdihedral} & CNOT-dihedral & width induction & 30 & 5{,}984 & sub-presentation, every $n$ \\
\citet{blake2026phaseaffine} & phase-affine, prime $d$ & width induction & 56 & 12{,}748 & sub-presentations, every $n$ \\
\citet{bian2023cliffordt2} & 2-qubit Clifford+$T$ & Reid.--Schreier & 62 & 12{,}371 & sub-presentation in $U_4$ \\
\citet{bian2023cliffordcs3} & 3-qubit Clifford+$CS$ & Reid.--Schreier & 123 & 57{,}744 & sub-presentation in $U_8$ \\
\citet{clement2026realcliffordch} & real Clifford+$CH$ & encoding & 346 & 74{,}422 & sub-presentation, every $n$ \\
\citet{amy2024toffoli} & 3-qubit Toffoli--Hadamard & encoding & 201 & 29{,}249 & presentation of $L_8$ \\
\citet{li2025qutrit} & qutrit Clifford & normal forms & 89 & 55{,}517 & presentation, every $n$ \\
\midrule
\multicolumn{6}{@{}l}{\emph{Minimality}}\\
\citet{blake2026simpler} & six fragments & models & 73 & 25{,}306 & minimality, completeness \\
\bottomrule
\end{tabularx}
\end{table}

\subsection{Two-level matrix groups}

Four of the sources present groups of two-level matrices by exact
synthesis, and the fifth (Hadamard-$\Pi$) does too; we formalized each
of them with the same skeleton, the levels and edges of
Section~\ref{sec:levels}. Greylyn's thesis~\cite{greylyn2014} presents
$U_4(\Z[\frac1{\sqrt2},i])$, the group of two-qubit Clifford+$T$
operators; our development states the relations for every $n$ and
proves completeness for $n\le4$, as the thesis does. Bian and
Selinger~\cite{bian2021un} present $U_n(\Z[\frac12,i])$, the group
behind Clifford+$CS$, for every $n$. Li, Ross and
Selinger~\cite{li2021on} present $O_n(\Z[\frac12])$ with the generators
$(-1)_{[a]}$, $X_{[a,b]}$ and $K_{[a,b,c,d]}=H\otimes H$, and the group
$L_n$ of scaled dyadic orthogonal matrices, which adds $I\otimes H$; we
formalized both, with normal syllables that carry one sign where theirs
carry four (Section~\ref{sec:beyond}).
Hadamard-$\Pi$~\cite{fang2026hadamard} presents $O_n(\Z[\frac1{\sqrt2}])$
by $Z_{[a]}$, $X_{[a,b]}$, $H_{[a,b]}$ and the relations (a1)--(d4) of
its Figure~6, and Cl\'ement's thesis of 2019~\cite{clement2019thesis}
the same group by relations (1)--(21); we formalized both theorems, the
first with the case its proof misses taken from the second
(Section~\ref{sec:gap}).

\subsection{Circuits}

Selinger's Clifford circuits~\cite{selinger2015clifford} and Makary,
Ross and Selinger's real stabilizer circuits~\cite{makary2021real} are
proved complete by normal forms, with supplements of derivations that
we followed line by line (Section~\ref{sec:search}). Both papers count
their normal forms, and so the order of the group; we proved the
counts as bijections with $\mathsf{Fin}$ of the order, and, for the
real case, that every orthogonal matrix normalizing the Pauli group
is the matrix of a circuit. Amy, Chen and Ross~\cite{amy2018cnotdihedral} present the
CNOT-dihedral operators, generated by $\omega$, $X$, $T$ and $CNOT$,
and Blake's phase-affine circuits~\cite{blake2026phaseaffine} form four
levels in each prime dimension $d$, affine circuits, then phases, then
$S$, then $T$, the last two for odd $d$ and $d>3$; both proofs go by
induction on the number of wires (Section~\ref{sec:width}), and the
known presentation of linear reversible circuits that the first builds
on is proved too, by a coset tower in the style of
Lafont~\cite{lafont2003boolean}. Bian and Selinger's two-qubit
Clifford+$T$ and three-qubit Clifford+$CS$
circuits~\cite{bian2023cliffordt2,bian2023cliffordcs3} came with Agda
code, which proves only that the circuit relations present the same
group as a set of two-level relations, along a translation of
generators, and does not mention matrices; we ported it and connected
it to the two-level presentations above at $n=4$ and $n=8$
(Section~\ref{sec:ports}), which gives the semantic theorem.
Cl\'ement's real Clifford+$CH$
circuits~\cite{clement2026realcliffordch} are proved complete width by
width: by a normal form on one qubit, and from two qubits on by
encoding circuits as words over two-level matrices, or over a group $P$
of pairs of them, and back, with the completeness of those words
imported from Hadamard-$\Pi$ (its Theorem~4.4) or proved by the
Reidemeister--Schreier method (its Theorem~4.10); its equations run to
(361), most of them auxiliary equations for the decoding, proved by
induction on the width with completeness one width down as a black
box. All of it is formalized (Section~\ref{sec:encoding}).

Two later developments complete the circuits.
\citet{amy2024toffoli} present the three-qubit Toffoli--Hadamard
circuits by the relations of their Figures~4, 7 and~8, through the
two-level presentations of \citet{li2021on}. We proved that these
relations, corrected and with one that the figures omit, present $L_8$,
the $8\times8$ scaled dyadic orthogonal matrices, by an encoding into
two-level words and back; without the omitted relation they do not
present $O_8(\Z[\frac12])$ (Section~\ref{sec:findings}).
\citet{li2025qutrit} give eighteen rules, C1--C18, for qutrit Clifford
circuits over $-\omega$, $H$, $S$ and $CZ$, complete by a normal form
of boxes. We proved that they present the $n$-qutrit Clifford group,
exact matrices over $\Z[\zeta_{12}]$, at every width, and that it has
order $6\prod_{k=1}^n 3(9^k-1)9^k$; the box relations that the normal
form needs are derived from the rules by generated derivations. The
Bian--Selinger and Toffoli--Hadamard theorems are also stated for
circuits in the library's circuit syntax.

\subsection{Qudit Clifford circuits}\label{sec:qupit}

\citet{bian2026qupit} give sixteen rules, C0--C15 of their Figure~1,
for Clifford circuits on qudits of odd prime dimension $p$, over
$-\omega$, $H$, $S$ and $CZ$. Their theorem was formalized alongside the
paper, with the same library, mainly by hand over about three months.
It goes up in layers, each a presentation theorem: $Sp(2n,\Z_p)$, by a
normal form computed by a coset tower; its semidirect product with the
Pauli group, the Clifford group modulo scalars
(Section~\ref{sec:products}); and the scalars, a central extension by
$\langle\omega\rangle\cong\Z_p$ given by the Weyl cocycle, and a direct
factor $\langle -1\rangle$. The group on the right is an algebraic
model, a Pauli and a symplectic matrix over $\Z_p$ with a phase and a
sign; that it is the group of unitary matrices the circuits generate
is the Weyl representation, standard but not formalized, so this
development says that the rules derive every equation that holds in the
model.

\subsection{Minimality and independence}\label{sec:structural}

Blake~\cite{blake2026simpler} gives smaller presentations for six
fragments of near-Clifford circuits and proves them minimal: no rule
follows from the others. The proofs use separating interpretations, a
model in which all rules but one hold. We formalized the six fragments
with their matrix semantics, soundness at every width, and all the
separating interpretations of the minimality theorem, on the
independence framework of Section~\ref{sec:independence-layer}. Since
the rules are sound, an independent rule is also necessary for
completeness, its own two sides being an equation that the semantics
validates and the other rules do not derive; the paper's Section~5.1
says so, and we state it for every fragment in the ranges of its
Theorem~37. The completeness part of that paper, a translation from
the earlier presentations, is formalized for all six fragments in the
ranges of its Theorem~22: at every width for the Clifford, real
Clifford, qutrit Clifford and CNOT-dihedral fragments, up to two wires
for Clifford+$T$ and three for Clifford+$CS$, from the presentations
above, with the derivations of its appendices generated by scripts.
So each rule set is a minimal complete presentation in those ranges;
and the Clifford+$CS$ rules, which derive the Clifford ones, are
complete for Clifford circuits at every width.
